\section{Model Training}
\subsection{Training Recipe}
Our training pipeline consists of three stages that progressively enhance model capabilities: foundational cross-modal alignment, multi-task pre-training, and supervised fine-tuning.
Table~\ref{tab:foundational-stages} summarizes the data composition, token scale, and training configurations for each stage.


\begin{table}[htbp]
\centering
\small 
\caption{Overview of the Training Stages. This includes details on data composition, token scale, sequence length, and trainable components.}
\label{tab:foundational-stages}
\renewcommand{\arraystretch}{1.0}
\resizebox{\linewidth}{!}{
\begin{tabular}{cccc}
\toprule
\textbf{Stages\&Objective} & \textbf{S0: Vision-Language Alignment} & \textbf{S1: Multi-task Pre-training} & \textbf{S2: Supervised Fine-Tuning}  \\
\midrule
\textbf{\makecell[c]{Understanding\\Data}} & \makecell[c]{Image Caption, Text} & \makecell[c]{Image Caption, Text, \\ OCR, Grounding, \\ Counting, Video Data,\\ Multimodal VQA} & \makecell[c]{High-quality Multimodal VQA\\ High-quality Text QA \\ Interleaved Reasoning} \\
\midrule
\textbf{\makecell[c]{Generation\\Data}} & Text-to-image & \makecell[c]{Text-to-image \\ Image Editing \\ Interleaved Generation}& \makecell[c]{High-quality Image Generation \\ Image Generation with CoT\\ High-quality Image Editing\\ High-quality Interleaved Generation\\ Interleaved Reasoning}  \\
\midrule
\textbf{Gen. Resolution}& $256\rightarrow512$ & 512& 512~~(diffusion decoder $\rightarrow$ 1024)  \\
\midrule
\textbf{Under. Max Edge}& 800 & 800& 800  \\
\midrule
\textbf{Training Tokens} & 100B & 210B & 80B \\
\midrule
\textbf{Sequence length} & 8192 & 8192& $8192\rightarrow16384$\\
\bottomrule
\end{tabular}
}
\vspace{-0.15in}
\end{table}

\noindent
\textbf{Stage 0: Vision–Language Alignment.} 
The primary objective of stage 0 (S0) is to align visual and linguistic representations within the dLLM backbone.
We use high-quality image–caption pairs and visual knowledge datasets, supplemented with pure text data to preserve language capabilities.
During training, a random masking strategy is applied to a subset of text and image tokens: for generation tasks, only image tokens are masked; for understanding tasks, only text tokens are masked.
To handle long visual token sequences, we adopt a progressive arbitrary resolution scheme: generation starts at $256 \times 256$ ($\sim$256 tokens) and transitions to $512 \times 512$ ($\sim$1024 tokens), while understanding consistently uses $800 \times 800$ with arbitrary resolution ($\sim$2048 tokens).



\noindent
\textbf{Stage 1: Multi-task Pre-training.} 
The model is trained on diverse multimodal data to develop comprehensive understanding and generation capabilities.
%
Visual understanding data includes image–text interleaved data, OCR, and visual counting/grounding tasks. 
Generation data includes image editing, subject-driven generation, controllable generation, style transfer using reference images, and multi-view generation tasks. 
%
This stage strengthens cross-modal connections and enables the model to handle increasingly complex tasks.

\noindent
\textbf{Stage 2: Supervised Fine-tuning.}
The Supervised Fine-Tuning (SFT) process is conducted in two stages: an initial phase at 8k context length for fundamental instruction-following capabilities, followed by expansion to 16k context for complex visual reasoning and generation.


\subsection{Pre-Training Optimization}
We adopt the Block Diffusion Language Model (BDLM)~\citep{arriola2025block} training objective, which extends standard discrete diffusion by operating on block-level masked regions rather than individual tokens.
%
This design enables parallel decoding while maintaining coherent context within each block, making it well-suited for the variable-length sequences common in multimodal tasks.

\noindent
\textbf{BDLM Loss.}
The training loss under the BDLM paradigm is defined as:
\begin{equation}\label{eq:bdlm}
    \mathcal{L}_{\text{BDLM}}(\theta) = - \mathbb{E}_{t, \bm{x}_0, \bm{x_t}} \left[ \frac{\alpha_{t}'}{1-\alpha_{t}}\sum_{k=1}^{K} \sum_{i=1}^{L_B} \mathbb{1}[x_{t,k}^i=\text{[MASK]}] \log p_{\theta}(\bm{x}_{0,k}^i | \bm{x}_{0,<k}, \bm{x}_{t,k}) \right],
\end{equation}
where the expectation is over timestep $t$, the clean sequence $\bm{x}_0$, and its corrupted version $\bm{x}_t$ (tokens masked with probability $1 - \alpha_t$). 
The indicator $\mathbb{1}[\cdot]$ ensures predictions are made only for masked tokens, and $-\alpha'_t / (1 - \alpha_t)$ is the diffusion-derived time weight. 
We define: $K = L_{\text{total}} / L_B$ as the number of blocks, $L_B$ as the block size, $x^i_{t,k}$ as the $i$-th token in block $k$, $\bm{x}_{0,<k}$ as the preceding clean blocks, and $\bm{x}_{t,k}$ as the noisy version of the current block.
\looseness=-1

\noindent
\textbf{Load Balancing Strategy.}
In MoE models, imbalanced expert utilization can lead to routing collapse.
%
We adopt an auxiliary-loss-free load balancing mechanism~\citep{deepseekv3} that promotes differentiated expert specialization while encouraging uniform workload distribution.
%
To improve numerical stability, we scale routing gate outputs by a factor of 2.5, stabilizing their root-mean-square (RMS) magnitude.
%
The auxiliary-loss-free bias is updated according to~\citep{sjl_moe}:
\begin{equation}
b_{i}=b_{i}+u\times \frac{(F_{i}-Q_{i})}{\sqrt{\frac{1}{n}\sum_{j=1}^{n}{(F_{j}-Q_{j})^2}}} \,,
\end{equation}
where $F=\mathbb{E}(f)$ denotes the current expert load distribution induced by the bias $b$, and $Q=[\frac{1}{n},\frac{1}{n},\dots,\frac{1}{n}]$ represents the ideal uniform distribution over $n$ experts. 
This RMSNorm-style normalization smooths bias updates, leading to stable load balancing throughout training.

\subsection{Supervised Fine-Tuning Optimization}
During SFT, we adopt the same load balancing strategy while introducing complementary masking and a mask token reweighting loss to handle variable-length sequences.

\noindent
\textbf{Mask Token Reweighting Loss.}
We adapt the BDLM objective to be conditional on an input prompt $c$:
\begin{equation}\label{eq:bdlm}
    \mathcal{L}_{\text{SFT}}(\theta) = - \mathbb{E}_{t, (c,\bm{x}_0), \bm{x_t}} \left[ \frac{\alpha_{t}'}{1-\alpha_{t}}\sum_{k=1}^{K} \sum_{i=1}^{L_B} \mathbb{1}[x_{t,k}^i=\text{[MASK]}] \log p_{\theta}(\bm{x}_{0,k}^i | c,\bm{x}_{0,<k}, \bm{x}_{t,k}) \right] \,.
\end{equation}

A key challenge in SFT is that sample lengths vary significantly---by up to two orders of magnitude.
Naive token-averaged loss causes gradients to be dominated by long sequences, while sample-level averaging encourages brevity.
We therefore propose a re-weighting mechanism to balance these extremes:
\begin{equation}
    \mathcal{L}_{\text{MTRS}} = \frac{\sum_{j} \beta_j \mathcal{L}_{\text{SFT}}^{(j)}}{\sum_{j} \beta_j} \,, \quad \text{where} \quad \beta_j = \frac{1}{\sqrt{\sum_{k=1}^{K} \sum_{i=1}^{L_B} \mathbb{1}[x_{t,k}^{i,(j)}=\text{[MASK]}]}} \,.
\end{equation}
The scaling factor $\beta_j$ is the inverse square root of the number of masked tokens in sample $j$, equilibrating gradient contributions across diverse response lengths.

\noindent
\textbf{Complementary Masking.}
Complementary masking~\citep{li2025lavida} is a strategy to enhance data efficiency for dLLMs by constructing two antithetical training instances from a single sequence $x_0$: a primary noised sequence $x_t$ and a complementary sequence $x'_t$ using the inverse mask.
This design ensures that every token position appears uncorrupted exactly once per pair, thereby doubling effective information utilization and eliminating token-level sampling bias.
We adapt this strategy in our framework.



\subsection{Diffusion Decoder Training}
\label{sec:decoder_training}
\paragraph{Training Paradigm.} 
We optimize the diffusion decoder via the standard flow matching objective~\citep{lipman2022flow}. The overall training trajectory is decoupled into a preliminary warm-up followed by a progressive two-stage fine-tuning scheme. The flow matching loss is formulated as:
\begin{equation}
    \mathcal{L}_{\text{FM}}(\theta) = \mathbb{E}_{\bm{x}_0, \bm{x}_1, \bm{z}, t} \left[ \| \bm{v}_{\theta,t}(\bm{x}_t, \bm{z}) - \bm{v}_{t} \|^2_2 \right],
\end{equation}
where $\bm{z}$ represents the conditioned semantic visual tokens, $\bm{v}_{\theta,t}$ denotes the velocity field predicted by the network at timestep $t$, and $\bm{v}_{t}$ is the target velocity. 
The overall training process is divided into three stages:
\begin{itemize}[leftmargin=16pt]
    \item \textbf{Stage 1: Warm-up.} We freeze the semantic processor and update only the remaining modules to establish cross-modal alignment while preserving pre-trained priors.

    \item \textbf{Stage 2: Multi-domain Generalization.} Following the warm-up, we unfreeze all parameters and
    fine-tune on diverse domains for robust generalization.

    \item \textbf{Stage 3: High-fidelity Refinement.} 
    In the final stage, we refine on high-quality data to elevate aesthetic fidelity and fine-grained visual details.
\end{itemize}

\paragraph{Few-step Generation.}
To accelerate visual generation, we adopt a lightweight consistency-based distillation framework~\citep{sun2026duality} for the diffusion decoder.
This method requires only an auxiliary projection layer---a final additional layer added to the decoder backbone---which will be discarded at inference time.
The distillation objective combines flow matching with a consistency term:
\looseness=-1
\begin{equation}
    \mathcal{L}_{\text{Distill}}(\theta) = 
    \mathbb{E}_{\bm{x}_0, \bm{z}, t} \left[ \| \bm{v}_{\theta,t} - \bm{v}_{t} \|^2_2 + \| \bm{u}_{\theta,t} - \bm{v}_{t} + t \cdot \frac{\text{d}\bm{u}_{\theta^-,t}}{\text{d}t} \|^2_2 \right], \text{where} \,\, \bm{u}_{\theta^-,t}=\text{stop\_grad(}\bm{u}_{\theta,t}\text{)} \,,
    \label{distill_loss}
\end{equation}
where $\bm{v}_{t}$ is the target velocity, $\bm{v}_{\theta,t}, \bm{u}_{\theta,t}$ are the dual outputs of the diffusion decoder.
% and the time derivative $\nicefrac{\text{d}\bm{u}_{\theta^-,t}}{\text{d}t}$ is a Jacobian-vector product (JVP)~\citep{lu2024simplifying} output of the diffusion decoder.
The time derivative $\nicefrac{\text{d}\bm{u}_{\theta^-,t}}{\text{d}t}$ is a Jacobian-vector product (JVP)~\citep{lu2024simplifying} output of the diffusion decoder, where the JVP calculation is approximated using the second-order difference technique proposed in UCGM~\citep{sun2025unified}.
This strategy enables 8-step CFG-free inference while maintaining high image quality.

\begin{figure}[htb]
    \centering
    \vspace{-0.15in}\includegraphics[width=0.85\linewidth]{figs/packing.pdf}
    \vspace{-0.05in}
    \caption{
    Data Packing Strategy for Efficient Training. Multiple shorter samples are concatenated into fixed-length sequences, minimizing padding tokens and improving GPU utilization.
    }
    \label{fig:pack}
    \vspace{-0.1in}
\end{figure}
\subsection{Infrastructure for Training Efficiency}
%
\noindent
\textbf{Image Tokens Pre-extraction.} 
LLaDA2.0-Uni employs a Vector Quantized (VQ) tokenizer to transform images into discrete visual tokens, incurring substantial computational cost during training.
%
We adopt an offline pre-extraction strategy: prior to training, the entire dataset is processed through the frozen tokenizer, with token indices stored on disk.
%
During training, the data loader retrieves pre-extracted tokens directly, eliminating repeated encoder passes and significantly accelerating the pipeline.

\noindent
\textbf{Load Balancing via Data Packing.}
In LLaDA2.0-Uni, sequence lengths vary significantly across multimodal tasks (e.g., short text-only tasks versus long image-generation tasks).
Traditional batching strategy requires extensive padding to match the longest sequence, wasting computation on padding tokens.
We address this with an offline data packing strategy that consolidates multiple shorter samples into fixed-length sequences (Figure~\ref{fig:pack}), significantly increasing effective token throughput for both pre-training and post-training.


\noindent
\textbf{Distributed Framework.}
We employ dFactory~\citep{dFactory} as the primary training engine for both pre-training and post-training phases---a high-efficiency framework specifically optimized for Diffusion Large Language Models. 
Built upon the VeOmni~\citep{ma2025veomni} distributed training ecosystem, dFactory enables the flexible deployment of sophisticated parallelization strategies. 